% Generado por Experimentos/generar_resultados_latex.py desde complementario.
% No editar a mano: volver a ejecutar el script.

\begin{table}[ht]
\centering
\footnotesize
\begin{tabular}{c c c c}
\hline
Ciclo & $\alpha=0$ & $\alpha=0.5$ & $\alpha=1$ \\
\hline
0 ($M_0$) & \multicolumn{3}{c}{$0.9133$ [$0.9111$; $0.9169$]} \\
1 & $0.9064$ [$0.9038$; $0.9077$] & $0.8981$ [$0.8958$; $0.8996$] & $0.8944$ [$0.8854$; $0.9057$] \\
2 & $0.9022$ [$0.8989$; $0.9050$] & $0.8962$ [$0.8793$; $0.9084$] & $0.8779$ [$0.8716$; $0.8874$] \\
3 & $0.9170$ [$0.9115$; $0.9226$] & $0.9000$ [$0.8900$; $0.9103$] & $0.8479$ [$0.8460$; $0.8498$] \\
4 & $0.9197$ [$0.9115$; $0.9310$] & $0.9017$ [$0.8908$; $0.9115$] & $0.8341$ [$0.8284$; $0.8372$] \\
5 & $0.9050$ [$0.8985$; $0.9142$] & $0.8959$ [$0.8900$; $0.9023$] & $0.8250$ [$0.8153$; $0.8349$] \\
\hline
\end{tabular}
\caption{Distinct-2: media y rango [mínimo; máximo] de las tres cadenas por ciclo y proporción sintética $\alpha$. $M_0$ es compartido por las tres proporciones dentro de cada repetición. Salvo la perplejidad, calculada con los 128 contextos de prueba, las métricas corresponden al análisis complementario \textit{post hoc} con los 90 contextos en los que todas las etapas alcanzan $W=30$ palabras.}
\label{tab:medias_distinct_2}
\end{table}
